\chapter{Notes}

\begin{table}[H]
\centering
\begin{tabular}{|p{3cm}|p{12cm}|}
\hline
\textbf{Test Case ID} & TC-01 \\ \hline
\textbf{Description} & Verification that the Makefile correctly performs the encryption step during the build process. The encrypted output file generated by the build system is compared with a previously generated reference file. The implementation details of the encryption process are referenced in the appendix. \\ \hline
\textbf{Prerequisites} & Build environment is set up and the Makefile encryption functionality is implemented. \\ \hline
\textbf{Test Steps} &
\begin{enumerate}
\item Build the project without the modified Makefile.
\item Rename the generated ELF file to keep it as a reference file.
\item Execute the build process again using the unchanged Makefile.
\item Compare the newly generated ELF file with the previously renamed reference file.
\item Verify that the files differ.
\end{enumerate}
\\ \hline
\textbf{Expected Result} & The encrypted output file generated during the second build shows changed areas to the previously generated reference file. \\ \hline
\textbf{Actual Result} & [Fill in] \\ \hline
\textbf{Status} & \textcolor{green!60!black}{Passed} or \textcolor{red}{Failed} \\ \hline
\end{tabular}
\end{table}

\begin{table}[H]
\centering
\begin{tabular}{|p{3cm}|p{12cm}|}
\hline
\textbf{Test Case ID} & TC-02 \\ \hline
\textbf{Description} & Verification that the Authenticator correctly transfers to the application after successful decryption and verification. The relevant implementation details are referenced in the appendix. [Fill in]\\ \hline
\textbf{Prerequisites} & Start Up of Authenticator and Application as well as copyAuthToRam, decrypt and verify functions as starting point for the test. \\ \hline
\textbf{Test Steps} &
\begin{enumerate}
\item Flash a simplified auth and app main that performs a simple LED indication.
\item Start the Authenticator.
\item Let the Authenticator copy the verify function to RAM.
\item Allow the Authenticator to perform decryption using a predefined key.
\item Verify the LED behaviour indicating that the application is running and the jump successful.
\end{enumerate}
\\ \hline
\textbf{Expected Result} & After successful verification, the Authenticator jumps to the application and the LED behaviour defined in the simplefied application becomes visible. \\ \hline
\textbf{Actual Result} & The jump to the application occurs and the switch in the LED behaviour confirms that the application code is executed. \\ \hline
\textbf{Status} & \textcolor{green!60!black}{Passed} \\ \hline
\end{tabular}
\end{table}

\begin{table}[H]
\centering
\begin{tabular}{|p{3cm}|p{12cm}|}
\hline
\textbf{Test Case ID} & TC-03 \\ \hline
\textbf{Description} & Verification that the implemented buffer overflow protection mechanisms prevent unintended behaviour when excessive input is received via the serial interface. The implementation details and scripts used for the test are referenced in the appendix. [Fill in] \\ \hline
\textbf{Prerequisites} & Full authenticator is operational. \\ \hline
\textbf{Test Steps} &
\begin{enumerate}
\item Start the full application and enter the initial key via the serial interface.
\item Verify correct system behaviour by placing a breakpoint at the subsequent state transition.
\item Execute the provided Python script to continuously send data to the serial interface.
\item Maintain the input stream for the duration of the timeout period (45 seconds).
\item Observe whether the system remains stable and switches to Failure.
\end{enumerate}
\\ \hline
\textbf{Expected Result} & The system remains stable during the entire input flooding period and it switches to Failure state after the given timeout period . \\ \hline
\textbf{Actual Result} & The authenticator continues running normally and the transition only occur after the expected timeout period. \\ \hline
\textbf{Status} & \textcolor{green!60!black}{Passed} \\ \hline
\end{tabular}
\end{table}

\begin{table}[H]
\centering
\begin{tabular}{|p{3cm}|p{12cm}|}
\hline
\textbf{Test Case ID} & TC-04 \\ \hline
\textbf{Description} & Verification of the functionality of the stack monitor module by observing stack usage before and after artificially increasing stack consumption. \\ \hline
\textbf{Prerequisites} & Stack is placed in linker and set up in start up; Stack monitor module is integrated and debugging tools are available. \\ \hline
\textbf{Test Steps} &
\begin{enumerate}
\item Execute all stack monitor module functions and store their initial values.
\item Observe the values using the debugger.
\item Introduce a large array in the code to intentionally increase stack usage (e.g. \textit{uint32\_t consumeStack[812]}).
\item Run the application again.
\item Observe the stack monitor values in the debugger.
\item Compare the values with the previously recorded measurements.
\end{enumerate}
\\ \hline
\textbf{Expected Result} & The stack monitor values change according to the increased stack usage and correctly reflect the consumed stack memory. \\ \hline
\textbf{Actual Result} & The monitored values indicate the expected change in stack usage after the additional memory consumption. \\ \hline
\textbf{Disclaimer} & This was tested on a different stack size than expected in the requirements and later implemented because the test took place before later changes in code. The viability of the test remains nonetheless. [Fill in reference] \\ \hline
\textbf{Status} & \textcolor{green!60!black}{Passed} \\ \hline
\end{tabular}
\end{table}

\begin{table}[H]
\centering
\begin{tabular}{|p{3cm}|p{12cm}|}
\hline
\textbf{Test Case ID} & TC-05 \\ \hline
\textbf{Description} & Verification that the stack monitor correctly detects stack corruption while the full state table of the application is active. \\ \hline
\textbf{Prerequisites} & Full application with active state table and stack monitor module is running. \\ \hline
\textbf{Test Steps} &
\begin{enumerate}
\item Run the application under normal operating conditions.
\item Monitor the return value of the \textit{isCorrupted()} function during execution.
\item Introduce a manual stack overflow by allocating an empty and a filled array on the stack.
\item Execute the application again with the modified code.
\item Observe whether \textit{isCorrupted()} returns true.
\end{enumerate}
\\ \hline
\textbf{Expected Result} & When stack corruption occurs, the \textit{isCorrupted()} function returns true and the system reacts accordingly. \\ \hline
\textbf{Actual Result} & The stack monitor correctly detects the corruption and reports the error condition when the array is filled. The stackoverflow by an array reservation alone is not detected.  \\ \hline
\textbf{Status} & \textcolor{orange!60!black}{Conditionally Passed} \\ \hline
\end{tabular}
\end{table}